\chapter{Melody and Counterpoint: Make One Sentence Memorable, Then Let the Second Answer It}
\label{ch:melody}

\classquote{凡音之起，由人心生也。人心之动，物使之然也。}{"Book of Rites", "Record of Music"}{All music begins in the human heart; when the heart is moved by things, it takes form in sound. Write the sentence you mean first; the melody comes after it. A melody is not better because it has more notes. First make one sentence stand, then decide whether the second repeats, answers or argues.}

\section{What a melody must do}
\label{sec:melody-job}

The melody layer does three things:

\begin{enumerate}
\item \textbf{Give identity}: the listener must recognise ``this is its sentence'', not a replaceable scale exercise.
\item \textbf{Give direction}: a sentence needs a beginning, a landing and a breath; staying in one register loses direction.
\item \textbf{Give emotion}: the same lyric feels different when the melody lands in major or minor, high or low, long or short phrases.
\end{enumerate}

A melody does not exist alone. It must land on the harmony's target notes, avoid the rhythm section's accents, and leave breath at section boundaries. Write the sentence first; write the notes second.

\section{Motif: short, memorable, changeable}
\label{sec:motif}

A \term{motif}{motif} is the shortest recognisable unit in a melody, usually 2--5 notes. A good motif is:

\begin{itemize}
\item \textbf{Short}: it can be said within two beats; longer and it becomes a phrase, hard to repeat.
\item \textbf{Characterful}: one leap, one syncopation, one special timbre or velocity makes it memorable.
\item \textbf{Changeable}: it can become a new sentence through transposition, inversion, shortening, extension or rhythmic change, rather than being copied exactly.
\end{itemize}

\begin{longtable}{p{0.20\textwidth}p{0.36\textwidth}p{0.36\textwidth}}
\toprule
Variation & How & Risk \\
\midrule
\endfirsthead
\toprule
Variation & How & Risk \\
\midrule
\endhead
Repetition & Play it exactly again & Reinforces memory; three times or more becomes tiresome. \\
Transposition & Move it up or down by several notes & Keeps the contour, changes colour; move too far and it leaves the harmony. \\
Inversion & Reverse the up/down direction & Creates an ``answer''; the listener may no longer recognise the original. \\
Shorten/extend & Cut or add notes at the ends & Changes the tone; cut too much and the motif loses its identity. \\
Rhythmic change & Keep the pitches, change durations/accents & Same motif, new character; change too much and it becomes another motif. \\
\bottomrule
\end{longtable}

Choose one feature when writing a motif; do not combine a leap, a syncopation and an ornament at once. More features make it harder to remember and harder to vary.

\section{Phrase: antecedent, consequent, breath}
\label{sec:phrase}

A \term{phrase}{phrase} is a melodic unit with a beginning and a landing, commonly 4 or 8 bars. Pop songs most often use an antecedent--consequent pair:

\begin{itemize}
\item \textbf{\term{Antecedent}{antecedent}}: poses the question and builds expectation, often ending on an unstable note or half cadence.
\item \textbf{\term{Consequent}{consequent}}: answers the question and resolves expectation, often returning to the tonic target note.
\item \textbf{Breath}: leave half a beat to a beat between phrases; without breath, the listener cannot hear where the sentence ends.
\end{itemize}

A practical check: sing every note of the melody. If you can only breathe after eight bars, the phrase is too long; if you break every two beats, it is too fragmented.

\section{Contour and target notes}
\label{sec:contour}

\term{Contour}{contour} is the melodic shape in pitch: rising, falling, arch, wave, staircase. It decides the listener's intuition about where the sentence is going.

\begin{itemize}
\item \textbf{Set the contour before filling notes}: a whole song circling one register feels flat; a whole phrase rushing upward leaves no room for the chorus.
\item \textbf{Target notes}: each phrase needs at least one \term{target note}{target-note} on the chord's root, third, fifth or seventh; the target decides whether the sentence is stable, suspended or resolved.
\item \textbf{Register management}: keep the verse's centre lower and lift it in the chorus; do not spend the highest note in the first phrase.
\item \textbf{Leaps need a fall}: after a large leap (five semitones or more), stepwise motion downward usually settles the line; it is not an absolute rule, but it is reliable.
\end{itemize}

\begin{figure}[htbp]
  \centering
  \begin{tikzpicture}[font=\small]
    \draw[->] (0,0) -- (9.2,0) node[right] {bars};
    \draw[->] (0,0) -- (0,3.8) node[above] {register};
    \draw[very thick, yebanblue] (0.4,1.2) -- (1.0,2.3) -- (1.8,1.8) -- (2.6,2.6) -- (3.4,2.0) -- (4.2,1.5) -- (5.2,1.3);
    \draw[thick, gray!70, dashed] (0.4,1.4) -- (1.4,1.8) -- (2.4,1.5) -- (3.4,2.2) -- (4.4,1.7) -- (5.2,1.9);
    \node[font=\scriptsize, align=center] at (0.7,0.5) {motif A};
    \node[font=\scriptsize, align=center] at (2.2,0.5) {A'};
    \node[font=\scriptsize, align=center] at (3.8,0.5) {contrast B};
    \node[font=\scriptsize, align=center] at (5.0,0.5) {return};
    \node[font=\scriptsize, anchor=west] at (6.0,2.4) {lead};
    \node[font=\scriptsize, anchor=west] at (6.0,1.8) {counter-line (dashed)};
  \end{tikzpicture}
  \caption{Motif--variation--contrast--return contour; the dashed counter-line often moves in the opposite direction so the two lines do not rise and fall together.}
  \label{fig:melody-contour}
\end{figure}

\section{Rhythm: accents, space and lyric}
\label{sec:melody-rhythm}

Melodic rhythm is not ``fill notes into the grid''; it is ``put the accent on the syllable that needs it'':

\begin{itemize}
\item \textbf{Lyric stress}: the stressed syllable of a multi-syllable word falls on a strong beat or a long note; do not put the accent on a function word.
\item \textbf{Syncopation}: syncopation creates drive, but constant syncopation loses the landing; each phrase needs one clear strong-beat landing.
\item \textbf{Space}: leave gaps between phrases and sections; silence is part of the arrangement.
\item \textbf{Density contrast}: the verse can be sparse and the chorus dense; do not use one density for the whole song.
\end{itemize}

\section{Counterpoint: how the second sentence answers the first}
\label{sec:counterpoint}

\term{Counterpoint}{counterpoint} is two or more melodic lines moving at once, independent yet coordinated. Pop rarely uses strict counterpoint; three practical relationships matter:

\begin{enumerate}
\item \textbf{Counter-melody}: write an independent line above or below the lead, in a different register, not rhythmically locked; do not copy the lead.
\item \textbf{Call-and-response}: after the lead phrase, another voice answers with a short phrase, usually 1--2 bars, without stealing the lead's identity.
\item \textbf{Contrary motion}: the lead rises while the counter-line falls, creating tension; do not keep both lines moving the same way forever.
\end{enumerate}

\begin{longtable}{p{0.22\textwidth}p{0.36\textwidth}p{0.34\textwidth}}
\toprule
Device & What it does & Risk \\
\midrule
\endfirsthead
\toprule
Device & What it does & Risk \\
\midrule
\endhead
Counter-melody & Different register and rhythm, independent phrase & Same register and rhythm as the lead turns it into a doubled melody. \\
Call-and-response & 1--2 bar answer after the lead & An answer that is too long or too loud steals the lead's identity. \\
Contrary motion & Lines move in opposite directions & Reversing for its own sake makes the harmony unstable. \\
\bottomrule
\end{longtable}

The bottom line of counterpoint is ``two roles must be audible'': the lead is the protagonist; the counter-line or answer is the supporting role. If both demand the lead's level, register and density, the listener hears one muddy line.

\section{Realising it in Yeban}
\label{sec:melody-yeban}

\begin{enumerate}
\item \textbf{Write a single line in the piano roll first.} Select a MIDI clip and press P / F7; write only the lead first, not the counter-line at the same time.
\item \textbf{Copy the clip for variations.} Use Ctrl+C / Ctrl+V to copy a motif clip, then change pitch, rhythm or length; do not rewrite the same motif by hand inside one clip.
\item \textbf{Transpose and check target notes.} Select notes and transpose the whole phrase; check that each phrase's target note lands on the current chord's root, third, fifth or seventh. Do not lift the whole part an octave just to sound ``high''.
\item \textbf{Use the velocity lane.} Draw velocity below the piano roll; separate the verse and chorus levels and leave breath inside phrases.
\item \textbf{Put the counter-line on its own track.} Do not stack it inside the lead track; create another MIDI track, keep the register apart, and solo both to check intelligibility.
\item \textbf{Check with Solo/Mute.} Solo the lead and confirm it stands alone; unmute the counter-line and confirm it does not cover the main phrase; finally unmute the rhythm section and confirm the drums do not eat the accents.
\end{enumerate}

\screenshot{arrangement-melody-counterpoint}{Lead and counter-line on two tracks in the piano roll, with the velocity lane: see separate registers, target notes and phrase breath}

\section{Common mistakes and fixes}
\label{sec:melody-mistakes}

\begin{longtable}{p{0.34\textwidth}p{0.58\textwidth}}
\toprule
Common mistake & Fix \\
\midrule
\endfirsthead
\toprule
Common mistake & Fix \\
\midrule
\endhead
Every note lands on a strong beat & Let weak beats and syncopation participate; keep one clear landing per phrase. \\
Too many notes in one phrase & Cut it down to something singable; use space and long notes for breath. \\
No recognisable motif & Choose 2--5 notes with one feature, then develop by repetition/variation. \\
The counter-line copies the lead & Separate register, rhythm and contour; the counter-line must stand as its own phrase. \\
Both lines always move the same way & Use contrary motion for tension, but serve the harmony rather than reversing mechanically. \\
The verse spends the highest note & Save the high note for the chorus; establish identity in the lower-middle register first. \\
Using transposition instead of melody writing & Transposition changes colour, not the sentence; change contour and rhythm first. \\
\bottomrule
\end{longtable}

\section{Exercise}
\label{sec:melody-exercise}

\begin{exercise}
Write an eight-bar lead and a four-bar counter-line:
\begin{enumerate}
\item Start with a three-note motif that can be said within two beats; then write variation A' (transpose or change rhythm) and a contrasting B.
\item Divide the eight bars into a four-bar antecedent and a four-bar consequent. End the antecedent on an unstable note, resolve the consequent to the tonic target note, and leave half a beat of breath between them.
\item Put the counter-line on a separate track: when the lead rises, the counter-line falls, and the registers stay apart. Solo both lines and confirm they are intelligible together.
\end{enumerate}
Finally, use the velocity lane to make a verse/chorus level difference and note which phrase's target note is clearest.
\end{exercise}

\section{Chapter takeaways}
\label{sec:melody-takeaways}

\begin{itemize}
\item A melody gives identity, direction and emotion before it worries about notes.
\item A motif must be short, memorable and changeable; vary one feature at a time.
\item Phrases have antecedent, consequent and breath; target notes land on chord tones.
\item Set the contour before filling notes; after a leap, stepwise fall usually settles the line.
\item Counterpoint must keep two roles audible: separate register and rhythm, and let the answer stay shorter than the lead.
\item Solve lead/counter-line conflicts at the arranging stage, not with volume.
\end{itemize}

\section{Further reading}
\label{sec:05-melody-further-reading}
\begin{itemize}
  \item \href{https://en.wikipedia.org/wiki/Melody}{Melody}: Revisit melody as a line of pitches and motives.
  \item \href{https://en.wikipedia.org/wiki/Songwriting}{Songwriting}: See how melody and lyric fit together.
  \item \href{https://www.youtube.com/watch?v=6gHEIF0rT2w}{Free music theory course}: Use scales and interval drills to support melodic writing.
\end{itemize}

\begin{helpmapping}
\begin{longtable}{p{0.45\textwidth}p{0.45\textwidth}}
\toprule
Section & Help anchor \\
\midrule
What a melody must do & \texttt{arrangement-melody-job} \\
Motif: short, memorable, changeable & \texttt{arrangement-motif} \\
Phrase: antecedent, consequent, breath & \texttt{arrangement-phrase} \\
Contour and target notes & \texttt{arrangement-contour} \\
Rhythm: accents, space and lyric & \texttt{arrangement-melody-rhythm} \\
Counterpoint & \texttt{arrangement-counterpoint} \\
Realising it in Yeban & \texttt{arrangement-melody-yeban} \\
Common mistakes and fixes & \texttt{arrangement-melody-mistakes} \\
Exercise & \texttt{arrangement-melody-exercise} \\
\bottomrule
\end{longtable}
\end{helpmapping}
